\section{Stable characters in every number of alphabets}
\label{sec:stable-characters}

The full character is the Euler character of the moment complex once
positive homology vanishes. The Schur resolution controls its earliest
possible degrees, and idempotent specialization excludes their boundary.
Quadratic exactness gives the complementary label bound. These two
inputs leave a short list of first-homology shapes; bivariate and hook
exactness remove the routine cases, and one explicit cubic multiplication
resolves the last one. The finer trace bounds are retained separately in
Appendix~\ref{app:finer-character-bounds}.

\subsection{Quadratic characters and exactness}
\label{sec:quadratic-characters}

Let $n\geq2$, $V=K^k$, and let $W_n$ carry its invariant
nondegenerate symmetric form. Removing the linear moments leaves
$P^0=\Sym(V\otimes W_n)$. The quadratic moment map identifies
$\Sym^2V$ with a subspace of $P^0_2$; put
\[
 Q_{n,k}=P^0/(\Sym^2V).
\]
The orthogonal Littlewood complex computes this quotient in every
rank.  Its modification rule also makes the quadratic Koszul complex
exact in positive homological degrees through total degree $n$.
These two facts give exact initial characters and provide the
homological input for the character range proved below.

For a graded $\GL(V)\times S_n$ character, let $\Pi_n$ retain just
the row Schur terms indexed by partitions $\rho$ with
$\rho'_1+\rho'_2\leq n-1$. Here primes denote conjugate partitions.

\begin{theorem}[Orthogonal quadratic character; Sam--Snowden--Weyman]
\label{thm:quadratic-character-all-ranks}
In the degree-completed representation ring,
\begin{equation}\label{eq:quadratic-projected-euler}
 [Q_{n,k}]=\Pi_n\left(
  [\Sym(V\otimes W_n)]
       \sum_{j\geq0}(-1)^j[\Lambda^j(\Sym^2V)]\right),
\end{equation}
where $\Lambda^j(\Sym^2V)$ has polynomial degree $2j$ and trivial
$S_n$-action. Equivalently,
\begin{equation}\label{eq:quadratic-orthogonal-character}
 Q_{n,k}\simeq
 \bigoplus_{\substack{\ell(\rho)\leq k\\
                       \rho'_1+\rho'_2\leq n-1}}
 \Schur_\rho V\otimes
       \left.\Schur_{[\rho]}W_n\right|_{S_n},
\end{equation}
where $\Schur_{[\rho]}$ is the orthogonal Schur module.
\end{theorem}
\begin{proof}
The Koszul complex of the quadratic moment subspace is
$\Sym(V\otimes W_n)\otimes\Lambda^\bullet(\Sym^2V)$.
Its row-$\rho$ multiplicity complex is the orthogonal Littlewood
complex.  The orthogonal Littlewood theorem identifies its
zeroth homology with $\Schur_{[\rho]}W_n$ for
$\rho'_1+\rho'_2\leq\dim W_n$, and zero otherwise; moreover all
positive homology vanishes for these admissible partitions
\cite[\S4.2, Proposition~4.1 and Theorem~4.4]{SSW2013}.
Taking Euler characteristics only on those row components proves
\eqref{eq:quadratic-projected-euler}; taking zeroth homology gives
\eqref{eq:quadratic-orthogonal-character}. Restriction from the full
orthogonal group to $S_n\subset O(W_n)$ retains all label
multiplicities. The construction and the identities are defined
over $\mathbb Q$, so characteristic-zero scalar extension gives
the assertion over $K$.
\end{proof}

\Needspace{12\baselineskip}
\begin{samepage}
\begin{corollary}[Three initial layers of every label]
\label{cor:first-three-character-layers}
Let $\mu\vdash n$, and let $\delta_\mu(k)$ be the least degree
in which $S^\mu$ occurs in the ambient polynomial ring $P_{n,k}$.
For a row partition $\rho\vdash D$ with $\ell(\rho)\le k$, the
joint $(\rho,\mu)$ multiplicities in $R_{n,k}[D]$ and $Q_{n,k}[D]$
agree whenever
\[
 D-\delta_\mu(k)\le2\qquad\text{or}\qquad\rho_1\le2.
\]
Here the empty partition has first part zero. Thus
\eqref{eq:quadratic-projected-euler} computes the complete character
in each of the three degrees
\[
 \delta_\mu(k),\quad\delta_\mu(k)+1,\quad\delta_\mu(k)+2,
\]
for every label and every number of alphabets, in every rank $n\ge2$.
It also computes the entire two-column row sector in all degrees.
For $n=1$, the quotient is $K$ and all nonconstant components vanish.
\end{corollary}
\end{samepage}
\begin{proof}
A polarized power sum of degree $d$ has row type $\Sym^dV$ and trivial label
type. To affect the joint $(\rho,\mu)$ component, it must multiply
a polynomial component of type $(\sigma,\mu)$ in degree $D-d$,
with $\rho/\sigma$ a horizontal $d$-strip. Consequently
\[
 d\le\rho_1,\qquad D-d\ge\delta_\mu(k).
\]
Either hypothesis of the corollary excludes every $d\ge3$. Thus higher
moments have zero image in the indicated component, even before
quotienting by the linear and quadratic moments. The quadratic
quotient therefore has the full coinvariant character there, as
given by Theorem~\ref{thm:quadratic-character-all-ranks}.
\end{proof}

Theorem~\ref{thm:r5-complete-first} computes $\delta_\mu(k)$ and
generates its entire first layer. The corollary computes the next
two full harmonic layers as well; these degrees may lie outside
the uniform stable range. Operator generation in the additional
layers is governed by the criteria in Section~\ref{sec:linear-generation}.

\begin{theorem}[Quadratic Koszul exactness through degree $n$]
\label{thm:quadratic-koszul-degree-range}
For every $n\ge2$, every finite-dimensional row space $V$, and
every $0\le D\le n$,
\begin{equation}\label{eq:quadratic-koszul-degree-range}
 H_j\!\left(\Sym(V\otimes W_n)
                 \otimes\Lambda^\bullet(\Sym^2V)\right)_D=0
                 \qquad(j>0),
\end{equation}
where the exterior generators have total degree two and the
differential is multiplication by the quadratic moments.
\end{theorem}

\begin{proof}
Put $m=n-1$ and decompose the complex by its row partitions
$\lambda\vdash D$.  Its row-$\lambda$ multiplicity complex is
the orthogonal Littlewood complex.  Write
$c_i=\lambda'_i$ for the column lengths.  If
$c_1+c_2\le m$, the admissible case of the Littlewood theorem
gives zero positive homology.

For an inadmissible partition of size at most $m+1$, the inequalities
\[
          m<c_1+c_2\le |\lambda|\le m+1
\]
force $c_1+c_2=m+1$.  In the orthogonal modification rule,
the dot action of the first reflection is
\[
 s_0\mathbin{\boldsymbol\cdot}(c_1,c_2,c_3,\ldots)
       =(m+1-c_2,m+1-c_1,c_3,\ldots).
\]
It fixes this column sequence.  The partition is therefore singular,
and its Littlewood complex has zero homology in every degree
\cite[\S4.4 and Theorem~4.4]{SSW2013}.
Every row partition of size $D\le n$ is either admissible or singular,
which proves \eqref{eq:quadratic-koszul-degree-range}.
\end{proof}

Consequently, in total degrees at most $n$, the projection $\Pi_n$
in \eqref{eq:quadratic-projected-euler} can be omitted: the ordinary
Koszul Euler character already equals $[Q_{n,k}]$.  Every discarded
row is singular and has Euler character zero in those degrees.
This exactness holds uniformly in the row dimension because it is
proved one Schur functor at a time; it requires no global
regular-sequence assertion for the quadratic moments.


\subsection{The full moment complex}
\label{sec:explicit-stable-characters}

Let $V=K^k$ and $P=P_{n,k}=\operatorname{Sym}(V\otimes K^n)$.
Write $p_\alpha=\sum_{i=1}^nX_i^\alpha$ for the polarized moment,
and give
$\operatorname{Sym}^j V$ weight $j$ in
$U_D=\bigoplus_{j=1}^D\operatorname{Sym}^j V$.
The moment Koszul complex is
\[
 \mathcal K_{n,k,D}=P\otimes\bigwedge U_D,
 \qquad \partial(e_\alpha)=p_\alpha.
\]
Here $e_\alpha$ is the exterior generator corresponding to the
monomial of row weight $\alpha$ in $\Sym^{|\alpha|}V$.
Its total grading counts both polynomial degree and the assigned
weights of exterior generators. Its zeroth homology agrees with
$R_{n,k}$ in total degrees at most $D$.

Our target is positive-degree exactness for
$D\ge3$, $n\ge\max\{D,2D-6\}$. Throughout this proof $D\le n$,
$\lambda$ is a row partition, and $\mu$ a label partition.
The Schur relation and quadratic reduction eliminate all homological
degrees at least two. At $n=2D-5$ the only remaining source has at
most two rows. At $n=2D-6$, three labels remain; their only component
not covered by bivariate or row-hook exactness is the multiplicity-one
cubic source in Lemma~\ref{lem:cubic-witness}.

\subsection{Parabolic detection of low-degree homology}

\begin{lemma}[Parabolic detection of low-degree homology]
\label{lem:parabolic-koszul-detection}
Fix $3\le D\le n$. In positive homological degree, every
label partition $\mu$ occurring in
$H_j(\mathcal K_{n,k,D})_D$ satisfies
\begin{equation}\label{eq:moment-koszul-homology-strip}
 n-D+3j\le\mu_1\le D-2.
\end{equation}
In any total degree, if the moments through order $a\ge1$ form a
regular sequence in $P$, one also has the lower bound
$\mu_1\ge n-D+(a+1)j$.
\end{lemma}

\begin{proof}
Remove the linear moments and compute quadratic Koszul homology first.
By Theorem~\ref{thm:quadratic-koszul-degree-range}, in total degrees at
most $n$ the remaining complex is
\[
 Q_{n,k}\otimes\Lambda^\bullet
       \left(\bigoplus_{d\ge3}\Sym^dV\right),
 \qquad Q_{n,k}=\Sym(V\otimes W_n)/(\Sym^2V).
\]
A degree-$b$ coefficient has first label row at least $n-b$, because
each monomial occupies at most $b$ labels; the bound passes to quotients.
The exterior factors consume at least $3j$ degrees, giving the lower
bound. If moments through order $a$ are regular, quotient by them first
and use the minimum remaining exterior weight $a+1$.

The trivial label has no positive moment homology through degree $n$,
by Theorem~\ref{thm:polynomial-schur-fronts}. For any other label, that
same theorem gives the stronger restriction
\begin{equation}\label{eq:polynomial-higher-label-support}
                  \mu_{j+1}\ge n-D+j+1.
\end{equation}
In particular $\mu_2\ge n-D+2$, so $\mu_1\le D-2$.
More generally, summing the first $j+1$ label rows yields
\[
 \mu_1\le jD-(j-1)n-j(j+1),\qquad
 n\le(1+1/j)D-j-4,\qquad D\ge j(j+4).
\]
The latter two inequalities combine this estimate with the quadratic
lower bound and $D\le n$. For example, second homology is absent in
this range whenever $D<12$.
\end{proof}


\begin{proposition}[Intersection support]
\label{prop:moment-intersection-support}
If $\Schur_\lambda V\otimes S^\mu$ occurs in $H_i$ in degree $D$
with $i\ge1$,
then, with $\rho_i,d_i$ from \eqref{eq:schur-resolution-shapes},
\[
 |\lambda\cap\rho_i|\ge2d_i-D,\qquad
 D+\lambda_1\ge n+\mu_1-\mu_{i+1}+i+1.
\]
Here the intersection is that of Young diagrams.
\end{proposition}
\begin{proof}
Inside $\Schur_{\rho_i}(A_+)_D$, at most $D-d_i$ inputs have degree
at least two. Splitting $A_+=V\oplus\Sym^{\ge2}V$ gives a
linear-input shape $\nu\subseteq\rho_i$ with $|\nu|\ge2d_i-D$.
Every resulting row contains $\nu$, proving the first bound; also
$\nu_1\ge(\rho_i)_1-(D-d_i)$ and $(\rho_i)_1=\mu_1+1$.
\end{proof}

Appendix~\ref{app:finer-character-bounds} gives further row-sensitive
trace bounds and detects their closed upper boundary. They are not
needed for the uniform character theorem below.

\subsection{Separation through degree \texorpdfstring{$n$}{n} in two alphabets}

A Hilbert-series identity alone need not eliminate cancelling homology.
Here multiplication already gives a free cover. Equality of dimensions
then forces its kernel to begin in a higher degree.

\begin{lemma}[Truncated freeness from a Hilbert identity]
\label{lem:truncated-freeness}
Let $S$ be a connected positively graded commutative or
supercommutative algebra over a field, and let $M$ be a nonnegatively
graded $S$-module. Assume their homogeneous pieces are finite
dimensional, and put $R=M/S_+M$. If
\[
 H_M(t)=H_S(t)H_R(t)\pmod{t^{N+1}},
\]
then every homogeneous section $R\to M$ induces an isomorphism
\[
 (S\otimes R)_{\le N}\longrightarrow M_{\le N}.
\]
Moreover, $\Tor_j^S(M,K)_d=0$ for $j>0$ and $d\le N$.
For superalgebras use ordinary dimensions and parity-preserving
sections. Additional positive multigradings may be retained.
\end{lemma}

\begin{proof}
The multiplication map is surjective by degree induction: after
subtracting a lift of the quotient class, an element belongs to
$S_+M$ and is a sum of positive-degree coefficients times
elements of smaller degree. The Hilbert identity makes the source
and target have equal dimensions through degree $N$. Its kernel
$L$ therefore has no homogeneous piece in those degrees.
Choose a graded free resolution of $L$ using generators only
in degrees greater than $N$, and splice it with
$0\to L\to S\otimes R\to M\to0$. Every positive term of the resulting
free resolution of $M$ is zero in degrees at most $N$.
Tensoring with $K$ proves the Tor assertion. All arguments respect
the additional gradings.
\end{proof}

\begin{theorem}[Bivariate separation through degree $n$]
\label{thm:bivariate-separation}
Multiplication induces a canonical $\GL_2\times S_n$-equivariant
isomorphism
\[
 \bigl(P_{n,2}^{S_n}\otimes\Harm_{n,2}\bigr)_{\le n}
          \xrightarrow{\ \sim\ }(P_{n,2})_{\le n}.
\]
Every positive moment Koszul homology group vanishes in internal
degrees at most $n$. Consequently the full bivariate $S_n$ character
is given, in bidegree $(a,b)$ whenever $a+b\le n$, by
\begin{equation}\label{eq:bivariate-character}
 \operatorname{tr}\bigl(\sigma\mid (R_{n,2})_{a,b}\bigr)
 =[q^at^b]\!
   \prod_{\substack{i,j\ge0\\i+j>0}}(1-q^it^j)
   \prod_{m\ge1}\bigl((1-q^m)(1-t^m)\bigr)^{-X_m(\sigma)}.
\end{equation}
Here $X_m(\sigma)$ counts the $m$-cycles of $\sigma$.
Equivalently, every row partition of length at most two has zero
positive moment homology through degree $n$, in every row dimension.
\end{theorem}

\begin{proof}
For every fixed bidegree $(a,b)$, with $D=a+b$, Wang's dimension
theorem \cite[Theorem~5.1]{Wang2026Polynomiality} gives one polynomial in $n$
equal to $\dim(R_{n,2})_{a,b}$ throughout $n\ge D$.
The Euler dimension of the moment complex is also polynomial in
$n$. Corollary~\ref{cor:universal-augmented-separation} gives eventual
acyclicity, for example for $n\ge2D-1$.
The actual and Euler dimension polynomials therefore agree at all
sufficiently large integers and are identical. Degrees zero,
one, and two follow from the independent linear and quadratic
moments. Hence the quotient has the Euler Hilbert series in every
total degree at most $n$.

Let
\[
 F_n=\Sym\left(\bigoplus_{j=1}^n\Sym^jK^2\right)
\]
act on $P_{n,2}$ by the formal moment evaluation. Polarized moments of
degrees at most $n$ generate the invariant algebra, so
$R_{n,2}=P_{n,2}/(F_n)_+P_{n,2}$.
The dimension
identity just proved is
$H_{P_{n,2}}=H_{F_n}H_{R_{n,2}}$ through degree $n$.
Lemma~\ref{lem:truncated-freeness} gives freeness and positive
Koszul acyclicity in this range.

The fundamental trace presentation has its first relation in
degree $n+1$: each relation uses $n+1$ positive inputs
\cite{Domokos2017}. Thus $F_n$ maps isomorphically
onto the actual invariant algebra in degrees at most $n$.
Choose the canonical harmonic section of $R_{n,2}$. The free
cover becomes exactly the multiplication map in the statement.

Taking equivariant characters of the moment complex gives
\eqref{eq:bivariate-character}. Each $m$-cycle contributes
$(1-q^m)^{-1}(1-t^m)^{-1}$ to the polynomial trace, and each
invariant moment contributes its Euler factor. Finally, the
complex is a polynomial functor of the row space. In
characteristic zero its homology decomposes into Schur functors;
every partition of length at most two is detected on $K^2$.
\end{proof}




\subsection{Row-hook exactness}
\label{sec:hook-homology}

The Fields character has a finite principal-specialization form
that records both the exactness range and its first obstruction.
Write $(q;q)_m=\prod_{i=1}^m(1-q^i)$ and
$(q;q)_\infty=\prod_{i\ge1}(1-q^i)$.
Let $Z$ be the Frobenius alphabet for $S_n$ and set
$X=Z/(1-q)$ in plethystic notation. This means
$p_d[X]=p_d[Z]/(1-q^d)$.

Let $\mathscr H_n(q,u;Z)$ be the Frobenius series of the
$(1,1)$ bosonic--fermionic coinvariant ring: $q$ records polynomial
degree and $u$ exterior degree. Let $\mathscr E_n$ be the
moment Euler character after evaluation on a row superspace of
dimension $(1|1)$.

\begin{lemma}[Finite Fields identity and its first discrepancy]
\label{lem:fields-euler}
For $0\le b\le n$,
\begin{equation}\label{eq:finite-fields}
 [u^b]\mathscr H_n
  =\sum_{h=0}^b(-1)^{b-h}
       \frac{(q;q)_{n-h}}{(q;q)_{b-h}}e_h[X]h_{n-h}[X].
\end{equation}
The corresponding coefficient of $\mathscr E_n$ is obtained by
replacing $(q;q)_{n-h}$ with $(q;q)_\infty$. Hence the two
characters agree in every total degree at most $n$, and
\begin{equation}\label{eq:first-fields-discrepancy}
 [q^{n-b+1}u^b](\mathscr H_n-\mathscr E_n)
       =e_b[Z]h_{n-b}[Z]\qquad(0\le b\le n).
\end{equation}
\end{lemma}

\begin{proof}
The Fields theorem \cite[Theorem~5.2 and equation~(2.5)]{MRW} gives
\[
 \mathscr H_n
 =\sum_{\mu\vdash n}s_\mu[Z]
      \sum_{T\in\SYT(\mu)}
        q^{\maj(T)}
        \prod_{r=1}^{\des(T)}(1+uq^{-r}).
\]
Put
\[
 \mathcal A_n(t,q;Z)=
   \sum_{\mu\vdash n}s_\mu[Z]
       \sum_{T\in\SYT(\mu)}t^{\des(T)}q^{\maj(T)}.
\]
Principal specialization of Schur functions gives the
character-valued Eulerian identity
\[
 \frac{\mathcal A_n(t,q;Z)}
      {\prod_{i=0}^n(1-tq^i)}
  =\sum_{m\ge0}
       h_n\left[Z\frac{1-q^{m+1}}{1-q}\right]t^m.
\]
One obtains it by standardizing semistandard tableaux and summing
the weakly increasing entries with forced rises at descents.
The standardization convention using comajor index gives the same
identity, since tableau evacuation preserves the number of
descents and interchanges major and comajor indices.

Expand
\[
 h_n[X-q^{m+1}X]
   =\sum_{h=0}^n(-1)^h q^{(m+1)h}e_h[X]h_{n-h}[X].
\]
Multiplication by the denominator above and evaluation at
$t=q^{-h}$ leave just one summand, yielding
\[
 \mathcal A_n(q^{-h},q;Z)
 =q^{-\binom h2}(q;q)_h(q;q)_{n-h}e_h[X]h_{n-h}[X].
\]
The elementary-symmetric identity
\[
 e_b(q^{-1},\ldots,q^{-d})
  =\sum_{h=0}^b
       \frac{(-1)^{b-h}q^{\binom h2}}
            {(q;q)_h(q;q)_{b-h}}q^{-hd}
\]
follows from the finite $q$-binomial theorem. Apply it with
$d=\des(T)$ and use the preceding evaluation to obtain
\eqref{eq:finite-fields}.

The polynomial--exterior algebra has Frobenius series
\[
 \sum_{h=0}^n u^h e_h[X]h_{n-h}[X].
\]
The formal moment algebra contributes the reciprocal Euler factor
\[
 \frac{(q;q)_\infty}{\prod_{a\ge0}(1+uq^a)},\qquad
 \prod_{a\ge0}(1+uq^a)^{-1}
     =\sum_{j\ge0}\frac{(-u)^j}{(q;q)_j}.
\]
This proves the asserted formula for $\mathscr E_n$.

The difference $(q;q)_{n-h}-(q;q)_\infty$ begins in degree
$n-h+1$ with coefficient one. In fixed exterior degree $b$,
every term of the difference therefore has $q$-degree at
least $n-b+1$. At that first degree only $h=b$ contributes,
and the constant term of $e_b[X]h_{n-b}[X]$ is
$e_b[Z]h_{n-b}[Z]$. This proves both conclusions.
\end{proof}

\begin{theorem}[Row-hook moment acyclicity]
\label{thm:hook-acyclicity}
For every $n$, every row dimension $k$, every row hook
$\lambda\vdash D$, and every $D\le n$,
\[
 \Hom_{\GL_k}\!\left(
   \Schur_\lambda K^k,H_i(\mathcal K_{n,k,D})_D\right)=0
                  \qquad(i>0).
\]
Thus the full $S_n$ character in every row-hook sector is its
moment Euler character throughout $D\le n$.
\end{theorem}

\begin{proof}
Evaluate the polynomial row functors on $K^{1|1}$. The coefficient
algebra is
$P=K[x_1,\ldots,x_n]\otimes\Lambda(\theta_1,\ldots,\theta_n)$.
The free super moment algebra $S$ has even generators of
bidegree $(a,0)$ for $a\ge1$ and odd generators of bidegree
$(a,1)$ for $a\ge0$. They act by
$\sum_i x_i^a$ and $\sum_i\theta_i x_i^a$, respectively.
Its ordinary Hilbert series is
\[
 H_S(q,u)=\frac{\prod_{a\ge0}(1+uq^a)}{(q;q)_\infty}.
\]
Trace generation identifies $P/S_+P$ with the superspace
coinvariant ring. Lemma~\ref{lem:fields-euler}, after taking
dimensions, says exactly that
$H_P=H_SH_{P/S_+P}$ through total degree $n$.
Lemma~\ref{lem:truncated-freeness} forces positive Tor to vanish
there. The super Koszul resolution uses exterior powers on even
generators and symmetric powers on odd generators, so it is
precisely this evaluation of the moment complex.

To pass back to ordinary row spaces, express each homogeneous
polynomial functor as
$\bigoplus_\lambda\Schur_\lambda\otimes C_\lambda$.
In characteristic zero all maps act on the multiplicity spaces
$C_\lambda$, and this decomposition can be evaluated on a
super vector space with its signed tensor permutations.
It commutes with homology. A nonempty Schur functor is nonzero
on $K^{1|1}$ exactly when its partition is a hook, and
\[
 s_{(a,1^b)}[q\mid u]=(q+u)q^{a-1}u^b.
\]
Vanishing of the specialized homology therefore forces the
vanishing of every row-hook multiplicity, including its entire
$S_n$-module structure.
\end{proof}



\subsection{Exact Euler characters and character polynomials}
\label{sec:exact-euler-characters}

\begin{theorem}[Stable character polynomial formula]
\label{thm:explicit-stable-character-polynomials}
For $D\ge3$ and $n\ge\max\{D,2D-6\}$, the total-degree-$D$ part of
the moment Koszul complex is exact in positive degrees.
The same conclusion holds for $D=0,1$ and $n\ge1$, and
for $D=2$ and $n\ge2$.
Consequently, if $\mathbf b=(b_1,\ldots,b_k)$ has total
size $D$, then for every $\sigma\in S_n$,
\begin{equation}\label{eq:explicit-stable-character-polynomial}
 \operatorname{tr}\bigl(\sigma\mid(R_{n,k})_{\mathbf b}\bigr)
 = [\mathbf q^{\mathbf b}]
   \prod_{0<|\alpha|\le D}(1-\mathbf q^\alpha)
   \prod_{j=1}^{D}\prod_{a=1}^{k}
                 (1-q_a^j)^{-X_j(\sigma)},
\end{equation}
where $X_j(\sigma)$ is the number of $j$-cycles of $\sigma$.
This gives the character polynomial explicitly in every
multidegree and every number of alphabets.

For $3\le D\le n$, the same Koszul product computes every
label coefficient $S^\mu$ with $\mu_1\ge D-1$.
\end{theorem}

\begin{proof}
Suppose $H_j$ is nonzero on label $\mu$ in degree $D\le n$.
Quadratic reduction gives $\mu_1\ge n-D+3j$.
Theorem~\ref{thm:polynomial-schur-fronts} excludes the trivial label
and gives $\mu_2\ge\mu_{j+1}\ge n-D+j+1$.
Adding the first two rows therefore gives $n\le2D-4j-1$.
Thus $j\ge2$ is impossible when $n\ge2D-6$.
For $j=1$ the same inequality gives $n\le2D-5$.
This proves exactness for $n\ge\max\{D,2D-4\}$.

At $n=2D-5\ge D$, the only possible label is $(D-2,D-3)$.
Its first homology is a subquotient of
$\Schur_{(D-1)}(A_+)_D=\Sym^{D-2}V\otimes\Sym^2V$ by the
relation presentation. This has at most two rows and is removed by
Theorem~\ref{thm:bivariate-separation}.

At $n=2D-6\ge D$, put $b=D-3\ge3$. The bounds
$\mu_1\ge b$, $\mu_2\ge b-1$ leave only
\[
               (b+1,b-1),\qquad(b,b),\qquad(b,b-1,1).
\]
For $(b+1,b-1)$ the first relation source has degree-$D$ part
$\Sym^{b+1}V\otimes\Sym^2V$, again with at most two rows.
The other two labels first occur in coefficient degree $b$, with
multiplicities $\Sym^bV$ and $\Schur_{(b-1,1)}V$, respectively.
Positive invariant relations cannot change this first occurrence.
Thus their degree-$(b+3)$ first chain groups after quadratic reduction
are exactly $\Sym^3V$ tensored with these spaces. The first has only
two rows.

For the remaining label $(b,b-1,1)$, the relation source is
$\Schur_{(b+1,1)}(A_+)_{b+3}$, whose one-quadratic-input decomposition is
\[
       \bigl(\Schur_{(b,1)}V\oplus\Sym^{b+1}V\bigr)
                                      \otimes\Sym^2V.
\]
Pieri leaves only $(b+1,1,1)$ and $(b,2,1)$ after the two-row
constituents are removed. Row-hook acyclicity
(Theorem~\ref{thm:hook-acyclicity}) removes the first. The second occurs
once in the cubic source
$\Sym^3V\otimes\Schur_{(b-1,1)}V$, and its differential is nonzero
by Lemma~\ref{lem:cubic-witness}. It is therefore injective there.
When $\dim V<3$, both three-row components are absent.

For $D=3,4,5$, $n\ge D$ leaves only the earlier cases; the cubic
witness is first needed at $D=n=6$. Degrees zero and one use the
linear sequence. In degree two the remaining map is the inclusion
of the independent quadratic moments for $n\ge2$.
The alternating character of the exterior factor is
$\prod_{0<|\alpha|\le D}(1-\mathbf q^\alpha)$, and a $j$-cycle
has polynomial trace $\prod_a(1-q_a^j)^{-1}$. Their product gives
\eqref{eq:explicit-stable-character-polynomial}.
The labelwise assertion follows from Lemma~\ref{lem:parabolic-koszul-detection}.
\end{proof}

\begin{corollary}[Separation of invariants and harmonics]
\label{cor:uniform-separation}
Multiplication induces an equivariant isomorphism
\[
 (P_{n,k}^{S_n}\otimes\Harm_{n,k})_{\le D}
                         \xrightarrow{\ \sim\ }(P_{n,k})_{\le D}
\]
when $k\le2$ and $D\le n$, or for arbitrary $k$ when
$D\ge3$ and $n\ge\max\{D,2D-6\}$.
For arbitrary $k$, its degree-$D$ part is also an isomorphism on
every row-hook isotype and every isotype whose row partition has
length at most two, whenever $D\le n$.
\end{corollary}
\begin{proof}
The harmonic inclusion is an equivariant section of the coinvariant
quotient. Apply Lemma~\ref{lem:truncated-freeness} to the character
identity, using the formal moment algebra, which agrees with
$P_{n,k}^{S_n}$ through degree $n$. The bivariate case is
Theorem~\ref{thm:bivariate-separation}; all smaller degrees satisfy
the same bounds, with degrees zero, one and two immediate.

For the componentwise assertion, let $L$ be the kernel of the free
cover over the formal moment algebra $S$. Reduction modulo $S_+$
is an isomorphism, so $L/S_+L\cong\Tor_1^S(P_{n,k},K)$.
A Littlewood--Richardson constituent contains the partition of
each input factor. A kernel component of least degree among row
hooks and partitions with at most two rows cannot lie in $S_+L$:
its $L$ factor would have lower degree and belong to one of those classes.
It would therefore survive in $\Tor_1$, contrary to
Theorems~\ref{thm:bivariate-separation} and~\ref{thm:hook-acyclicity}
through degree $n$.
\end{proof}

The coefficient extraction in
\eqref{eq:explicit-stable-character-polynomial} is finite.
Expanding $(1-z)^{-X}$ as
$\sum_{v\ge0}\binom{X+v-1}{v}z^v$ shows directly that
it is a rational polynomial in the cycle counts, of weighted
degree at most $D$ when $\deg X_j=j$. At the identity it gives
\begin{equation}\label{eq:explicit-stable-dimension-polynomial}
 \dim(R_{n,k})_{\mathbf b}
 = [\mathbf q^{\mathbf b}]
    \prod_{0<|\alpha|\le D}(1-\mathbf q^\alpha)
    \prod_{a=1}^k(1-q_a)^{-n}.
\end{equation}
Its degree in $n$ is exactly $D$, with leading coefficient
$1/(b_1!\cdots b_k!)$.
For total degree $D$, the corresponding dimension is
\[
 [t^D](1-t)^{-k(n-1)}
       \prod_{j=2}^{D}(1-t^j)^{\binom{k+j-1}{j}},
\]
with leading term $k^Dn^D/D!$.
For $D\ge3$ and $n\ge2D-3$, Theorem~\ref{thm:linear-generation}
gives these same characters for the Vandermonde operator module.

This explicitly describes the eventual character and dimension
polynomials asked for in
\cite[Problems~1.10 and~1.12]{ChurchEllenbergFarb2015}.
The existence of eventual polynomiality and representation
stability is due to Church--Ellenberg--Farb; see also the
quantitative refinements of \cite{Bahran2024,Borisov2026}.
Outside this range, the same bounds locate every possible homological
correction.  The Euler product is the low-degree quotient proposed in
\cite[Section~4, equation~(33)]{Bergeron2013}.
Bergeron's conjecture asks for its validity throughout $D\le n$.


\subsection{Two-column homology, reflected cycles, and multiplication}
\label{sec:two-column-homology}

Put $W=W_n$ and $m=\dim W=n-1$. For $0\le q\le p$, let
$\lambda(p,q)=(2^q,1^{p-q})$, the partition with column heights $p,q$.
Its degree is $D=p+q$. Write $M_{p,q}$ for its stable row multiplicity
space in full diagonal harmonics, with $M_{0,0}=K$.
Only the invariant quadratic form survives in this sector; consequently
its multiplicity spaces and moment homology carry an $O(W)$ action
extending their $S_n$ action.

\begin{theorem}[Two-term reduction in every degree]
\label{thm:two-column-complex}
After eliminating the linear moments, the full moment Koszul complex
in row $\lambda(p,q)$ is
\begin{equation}\label{eq:two-column-two-term}
 \Schur_{\lambda(p-1,q-1)}W
                  \xrightarrow{\ L\ }
 \Schur_{\lambda(p,q)}W
\end{equation}
in homological degrees one and zero, with zero source when $q=0$.
The differential inserts the invariant symmetric form.
There is no $H_j$ for $j\ge2$ in any two-column row.
\end{theorem}
\begin{proof}
A moment of order at least three forces first row at least three
in every Littlewood--Richardson constituent. The same is true of
two or more quadratic Koszul generators, since
$\bigwedge^2\Sym^2V=\Schur_{(3,1)}V$ and
$\bigwedge^j\Sym^2V$ is a summand of
$\bigwedge^2\Sym^2V\otimes\bigwedge^{j-2}\Sym^2V$ for $j\ge2$.
Thus only zero or one quadratic generator survives. Cauchy
decomposition gives the target; horizontal Pieri removes one box
from each column to give the source once. The differential is
quadratic multiplication.
\end{proof}

Let $U=K^2$ record two exterior copies of $W$, and put
\[
 \mathcal E=\Lambda(U\otimes W),\qquad
 \omega=\sum_{i=1}^m\theta_i\wedge\phi_i,\qquad
 L=\omega\wedge-,\qquad
 c=\sum_{i=1}^m\partial_{\phi_i}\partial_{\theta_i}.
\]
The displayed orthonormal notation denotes the operators defined by
the form and its inverse, so no choice of an orthonormal basis over
$K$ is required. They satisfy
$[c,L]|_{\mathcal E_e}=(m-e)\operatorname{id}$. For
$P_e=\ker(c:\mathcal E_e\to\mathcal E_{e-2})$, $0\le e\le m$,
the exterior Lefschetz decomposition gives
\begin{align}
 \mathcal E&=\bigoplus_{e=0}^m\ \bigoplus_{j=0}^{m-e}L^jP_e,
 \label{eq:exterior-lefschetz-strings}\\
 cL^jf&=j(m-e-j+1)L^{j-1}f\qquad(f\in P_e).
 \label{eq:exterior-string-normalization}
\end{align}
These follow directly from finite-dimensional $\mathfrak{sl}_2$
theory; see also the exterior Lefschetz treatment in
\cite{KimRhoades}. The operator $L$ commutes with the row
$\mathfrak{sl}_2$ on $U$. On bidegree $(a,b)$, its formula in
$\Lambda^aW\otimes\Lambda^bW$ is
$(-1)^a\sum_i\eps_i\otimes\eps_i$. Exterior multiplication by
$\omega$ fixes these signs uniformly.

Write $\eta$ for the quadratic Koszul generator, of homological degree
one, internal degree two, and $GL(U)$ type $\det U$. The reduced
complex is $\eta\mathcal E\xrightarrow{\,L\,}\mathcal E$, with
$d(\eta f)=Lf$. For its algebra formulas we use the totalized
product, so $\eta$ anticommutes with $U\otimes W$; the complete
sign convention is given in Appendix~\ref{app:odd-derived}.

\begin{theorem}[All two-column homology and its reflection]
\label{thm:two-column-reflection}
For $0\le q\le p\le n$, set $(p^*,q^*)=(n-q,n-p)$.
\begin{enumerate}
\item If $p+q\le n-1$, the sole homology is
      $H_0=M_{p,q}=\Schur_{[\lambda(p,q)]}W$ as $O(W)$ modules.
\item If $p+q=n$, the complex is acyclic.
\item If $p+q\ge n+1$, the sole homology is
      \[
       H_1\cong M_{p^*,q^*}
          \cong M_{p^*,q^*}^{\,*}\otimes(\det W)^2
      \]
      as $O(W)$ modules, hence as $S_n$ modules.
      The reflected row has degree $2n-(p+q)$.
\end{enumerate}
If $p>n$, both terms vanish. These assertions hold in every row
dimension in which the original row occurs; in case~(3),
$p^*\le p-1$, so the reflected row occurs there as well.

More explicitly, if $f\in P_e$ and $k=m-e$, reflection is the cycle map
\begin{equation}\label{eq:exterior-reflected-cycle}
 \rho_e(f)=\frac{\eta L^kf}{k!},
 \qquad
 \rho_e^{-1}(\eta g)=\frac{c^kg}{k!}.
\end{equation}
Before extracting row multiplicity spaces, its source is
$P_e\otimes(\det U)^{k+1}$ and its target has internal degree $2n-e$.
\end{theorem}
\begin{proof}
Each string in \eqref{eq:exterior-lefschetz-strings} contributes its
initial vector to $\mathcal E/L\mathcal E$ and its terminal vector to
$\ker L$. Equation~\eqref{eq:exterior-string-normalization} gives
$L^{k+1}f=0$ and $c^kL^kf=(k!)^2f$, proving
\eqref{eq:exterior-reflected-cycle}.

The row-$\lambda(p,q)$ multiplicity space is the highest-weight kernel
of the transfer
$\Lambda^pW\otimes\Lambda^qW\to
\Lambda^{p+1}W\otimes\Lambda^{q-1}W$.
These kernels fix the Pieri normalization of
\eqref{eq:two-column-two-term}. Its low-degree primitive part is the
orthogonal traceless Schur module: within-column traces vanish by
antisymmetry, and the remaining trace is encoded by $c$.
This is also the admissible case of the orthogonal Littlewood theorem
\cite[Theorem~4.4]{SSW2013}.
If $f$ has $GL(U)$ highest weight $(a,b)$, then $e=a+b$ and
\eqref{eq:exterior-reflected-cycle} has weight
\[
       (a+k+1,b+k+1)=(n-b,n-a).
\]
Exterior Cauchy decomposition transposes this to the stated row
reflection. The strings leave no homology in internal degree $n$.
The dual formula follows from the exterior top-product pairing,
whose orientation after adjoining $\eta$ is
$(\det U)^n\otimes(\det W)^2$.
Finally, if $p>n$, both Schur terms have more than $\dim W$ rows.
\end{proof}

For the bivariate $S_n$-character $R_n(u,v)$ of the exterior $H_0$,
reflection and the Euler character give
\[
 R_n(u,v)-(uv)^nR_n(u^{-1},v^{-1})
    =(1-uv)\lambda_u(W_n)\lambda_v(W_n),
 \qquad
 \lambda_u(W_n)=\sum_j[\Lambda^jW_n]u^j.
\]
Products are tensor products in the representation ring. The dual
operation disappears because $S_n$-modules are self-dual.

\begin{corollary}[The full two-column first syzygies]
\label{cor:two-column-first-syzygies}
In degree $n+1$, for $1\le b\le\lfloor(n+1)/2\rfloor$, the row
$\lambda=(2^b,1^{n+1-2b})$ has
\[
 H_{1,\lambda}
 \cong\Schur_{[(2^{b-1},1^{n+1-2b})]}W_n\big|_{S_n}
 \cong\sgn\otimes\Schur_{[(2^{b-1})]}W_n\big|_{S_n}.
\]
In particular, for $n\ge4$,
\[
 H_{1,(2,2,1^{n-3})}
   \cong S^{(2,1^{n-2})}\oplus S^{(2,2,1^{n-4})}.
\]
\end{corollary}
\begin{proof}
At this boundary the reflected columns are obtained by subtracting
one from each original column. Orthogonal determinant duality
complements the first column in dimension $n-1$, giving the second
formula. For $b=2$ use
$\Schur_{[(2)]}W_n=\Sym^2W_n-K
 \cong S^{(n-1,1)}\oplus S^{(n-2,2)}$ for $n\ge4$.
\end{proof}

Thus at $n=3$ row $(2,2)$ reflects to $(2)$ and has
$H_1\cong S^{(2,1)}$; at $n=4$ row $(2,2,1)$ reflects to $(2,1)$
and has $H_1\cong S^{(2,1,1)}\oplus S^{(2,2)}$.
The exterior $H_0$ character is also obtained in
\cite{KimRhoades,Lentfer}, with explicit double-hook label formulas
in \cite{LentferSign}; the positive homology is the width-two case
of \cite[Theorem~4.4]{SSW2013}. The cycle map above will now identify
its multiplication as well.

\begin{proposition}[Multiplication on reflected homology]
\label{prop:two-column-module}
Put $R=\mathcal E/(\omega)$ and
$I=\eta\operatorname{Ann}_{\mathcal E}(\omega)$.
For $\beta\in U\otimes W$, set
$\iota_\beta=[c,\varepsilon_\beta]$. Then
\begin{equation}\label{eq:two-column-contraction-action}
 \beta\,\rho_e(f)=-\rho_{e-1}(\iota_\beta f),
 \qquad
 I=R\cdot(\eta P_m).
\end{equation}
In particular, the minimal $R$-generator space of $I$ is $\eta P_m$,
all in internal degree $n+1$, and
\[
 \dim(I/R_+I)
   =\binom{2m}{m}-\binom{2m}{m-2}
   =\frac{1}{n+1}\binom{2n}{n}.
\]
The last formula uses the usual zero convention for binomial
coefficients outside their range.
\end{proposition}
\begin{proof}
The identities $[L,\iota_\beta]=\varepsilon_\beta$ and
$[c,\iota_\beta]=0$ give
$\iota_\beta P_e\subseteq P_{e-1}$. With $k=m-e$,
\[
 L^{k+1}\iota_\beta f=(k+1)\beta L^kf
\]
because $L^{k+1}f=0$. Division by $(k+1)!$, followed by moving
$\beta$ across $\eta$, gives the first identity.
The contraction $U\otimes W\otimes P_e\to P_{e-1}$ is a nonzero
symplectic-equivariant map onto an irreducible primitive module.
Contracting $\theta_1\cdots\theta_e$ by $\phi_e$ proves nonvanishing.
It is therefore onto, and iteration from $e=m$ proves generation.
For $m=0$ the assertion is simply $I=K\eta$.
The generator quotient is its first nonzero degree; the dimension
formula is the primitive exterior-power formula.
\end{proof}

The Catalan formula of \cite[Corollary~4.4]{KimRhoades}
counts the minimal syzygy-generating layer after evaluation on two odd rows.
The identities are $Sp_{2m}$-equivariant. For $GL(U)\times O(W)$,
contraction contributes the extra $\det U$ needed between the source
twists of $\rho_e$ and $\rho_{e-1}$; on row multiplicity spaces they
are $O(W)$-equivariant.

\begin{corollary}[All two-column syzygies are generated at the first boundary]
\label{cor:two-column-syzygy-generation}
Write $H_{j,D}(V)$ for the degree-$D$ moment homology as a polynomial
functor of the ordinary row space $V$. For every $D\ge n+1$, the
cokernel of multiplication
\[
 H_{0,D-n-1}(V)\otimes H_{1,n+1}(V)
                         \longrightarrow H_{1,D}(V)
\]
has no constituent $\Schur_\lambda V$ with $\lambda_1\le2$.
This holds in every row dimension in which that constituent occurs.
\end{corollary}
\begin{proof}
Purely odd two-dimensional evaluation is exact in characteristic
zero and sends
$\Schur_\lambda(\Pi U)$ to
$\Pi^{|\lambda|}\Schur_{\lambda'}U$.
It kills precisely the shapes with $\lambda_1>2$ and detects every
remaining Schur multiplicity. It preserves the moment differential
and multiplication. The evaluated map is onto by
\eqref{eq:two-column-contraction-action}, so its original cokernel
has no surviving two-column shape.
\end{proof}

The quotient by shapes with more than two columns carries the induced
algebra structure, since these shapes form a tensor ideal by
Littlewood--Richardson containment.
Appendix~\ref{app:odd-derived} identifies the complete odd differential
algebra and shows that triple Massey products generate its positive
homology.

\begin{proposition}[A single-degree criterion for two-column generation]
\label{prop:two-column-middle-test}
For fixed $n\ge2$, Vandermonde generation in every two-column row and
every degree is equivalent to generation of the two-column rows in
degree $n-1$. It suffices to test the rows
\[
             (2^q,1^{n-1-2q}),
       \qquad 2\le q\le\lfloor(n-1)/2\rfloor.
\]
For $n\ge5$, these tests require only $n-3$ ordinary alphabets.
All two-column rows are generated for $n\le5$.
\end{proposition}
\begin{proof}
Let $G$ be the purely odd two-dimensional evaluation of the ordinary
operator subfunctor $\Pol$. This means evaluation of the entire
operator module, not of its Vandermonde seed alone.
The corresponding full harmonic space is
$\bigoplus_{e=0}^mP_e$.
Derivative closure is expressed by the natural polynomial coaction
\[
 \Pol_D(V)\longrightarrow
            (V\otimes W)\otimes\Pol_{D-1}(V).
\]
Exact odd evaluation therefore makes $G$ closed under every
constant-coefficient exterior contraction.
The map $(U\otimes W)^*\otimes P_e\to P_{e-1}$ is onto, by the
same nonzero symplectic-equivariant contraction used in
Proposition~\ref{prop:two-column-module}.
Thus $G_m=P_m$ propagates down to $G_e=P_e$ for every $e\le m$.
Odd evaluation detects every two-column shape, proving the
equivalence. The cases $q=0,1$ are row hooks and follow from
Proposition~\ref{prop:row-hook-generation}. Each remaining row has
$n-1-q\le n-3$ parts, so Lemma~\ref{lem:stability} and
Corollary~\ref{cor:r3-exact-specialization} make $n-3$ alphabets
detect every remaining multiplicity. The list is empty for $n\le4$;
for $n=5$ its only row is $(2,2)$, covered by Haiman's bivariate
operator theorem as in Theorem~\ref{thm:long-first-row}.
\end{proof}

The remaining task is to generate the displayed degree-$(n-1)$ rows.
The derivative coaction then propagates generation to every lower degree.

\subsection{The first row-hook moment syzygies}

\begin{theorem}[The first row-hook moment syzygies]
\label{thm:first-hook-syzygies}
In total degree $n+1$,
\[
 \left.H_1(\mathcal K_{n,k,n+1})_{n+1}
                 \right|_{\mathrm{row\ hooks}}
  \cong
  \bigoplus_{b=0}^{n-1}
      \Schur_{(n+1-b,1^b)}K^k\otimes S^{(n-b,1^b)}.
\]
Each summand occurs once. Row lengths exceeding $k$ give zero
Schur modules. Every $H_i$ with $i\ge2$ has zero row-hook part
in this total degree.
\end{theorem}

\begin{proof}
In the $(1|1)$ free cover from the proof of Theorem~\ref{thm:hook-acyclicity}, its kernel
begins in degree $n+1$. A positive term of the super Koszul
complex on that kernel has internal degree greater than
$n+1$. The Tor long exact sequence therefore makes every
$H_i$, $i\ge2$, zero in degree $n+1$.
Consequently the character discrepancy
\eqref{eq:first-fields-discrepancy} is exactly $H_1$.

Let $T_b$ be the label multiplicity space of the row
$(n+1-b,1^b)$ in this homology, and put $T_{-1}=0$.
The two monomials in its super Schur specialization give
\[
 \operatorname{ch}T_b+\operatorname{ch}T_{b-1}
       =e_b[Z]h_{n-b}[Z]\qquad(0\le b\le n).
\]
For $1\le b\le n-1$, Pieri's rule says
\[
 e_bh_{n-b}
    =s_{(n-b,1^b)}+s_{(n-b+1,1^{b-1})}.
\]
At the endpoints the right side is $h_n$ and $e_n$.
Induction gives $T_b\cong S^{(n-b,1^b)}$ for $0\le b\le n-1$,
and $T_n=0$. Exact Schur evaluation yields the stated
ordinary-row result and the higher-homology vanishing.
\end{proof}

For $n=1$, the theorem gives only $\Schur_{(2)}K^k$ with
trivial label in degree two. For $n=2$, it gives
$\Schur_{(3)}K^k\otimes S^{(2)}$ and
$\Schur_{(2,1)}K^k\otimes S^{(1,1)}$ in degree three.
These also follow by direct moment calculation.



\subsection{Superspace characters and the first correction}

Theorem~\ref{thm:first-unstable-specht-wall} identifies the first
unstable correction for every nonempty tail $\tau$ and every
$D>|\tau|$. Its idempotent specialization proves injectivity at
that wall directly, without an Euler-character comparison.

\begin{corollary}[Stable diagonal superspace characters]
\label{cor:stable-superspace-characters}
Let $H_+(X)=\sum_{d\ge1}h_d(X)$ and define
$\operatorname{sp}_{Q\mid U}(p_m)=p_m(Q)+(-1)^{m-1}p_m(U)$.
For $\tau\vdash s$, $D\ge s$ and $n\ge D+\tau_1$, the
degree-$D$ multiplicity character in any bosonic--fermionic superspace is
\[
 \operatorname{ch}\operatorname{Hom}_{S_n}
 (S^{\tau[n]},R_n^{(k,j)}[D])
 =\operatorname{sp}_{Q\mid U}\bigl([s_\tau[H_+(X)]]_D\bigr).
\]
The full character-polynomial formula of
Theorem~\ref{thm:explicit-stable-character-polynomials} likewise
transfers to arbitrary superspaces in its range
$D\ge3$, $n\ge\max\{D,2D-6\}$.
\end{corollary}
\begin{proof}
The bosonic identities hold in every row dimension and hence are
identities of universal row symmetric functions. Apply Lentfer's
diagonal supersymmetry theorem~\cite[Theorem~1.2 and Corollary~1.3]{Lentfer}.
The displayed specialization records parity; at module level the
corresponding functors are super Schur functors.
\end{proof}
